\documentclass[../../main.tex]{subfiles}
\begin{document}

{\bf [解]}
(1)
$_{l+n}C_l = \dfrac{(l+n)!}{l!\cdot n!}$ だから与えられた条件は
\begin{align}
  {}_{l+n}C_l       &> {}_{m+n}C_m \\
  \dfrac{(l+n)!}{l!} &> \dfrac{(m+n)!}{m!} \\
  (l+n)\cdots(l+1)  &> (m+n)\cdots(m+1) \\
   \therefore
  l&>m
\end{align}
と変形できる．よって求める条件は
\begin{align}
 l > m
\end{align}
である．

(2)
$R(x)$ は高々2次式だから
\begin{align}
 R(x)=A(x-a)^2+B(x-a)+C \label{eq:}
\end{align}
とおける．
また適当な多項式 $P(x)$ があって，
\begin{align}
  F(x)=(x-a)^3P(x)+R(x) \label{eq:2}
\end{align}
とかける．
ここで$G(x)=F(x)-R(x)$ とおくと\cref{eq:2}より
\begin{align}
 G(x) = (x-a)^3P(x)
\end{align}
とかけて$G(x)$ は $(x-a)^3$ で割り切れるから
\begin{align}
  & G(a)=G'(a)=G''(a)=0 \\
\therefore
  & R(a)=F(a),\ R'(a)=F'(a),\ R''(a)=F''(a) \quad \cdots \text{①} \label{eq:3}
\end{align}
と書ける．

一方で$R$の微分は\cref{eq:1}より
$R(a)=C,\ R'(a)=B,\ R''(a)=2A$ だから，\cref{eq:3}より
\begin{align}
  C=F(a),\ B=F'(a),\ A=\dfrac{1}{2}F''(a)
\end{align}
と書けるから\cref{eq:1}に代入して
\begin{align}
  R(x)=\dfrac{1}{2}F''(a)(x-a)^2+F'(a)(x-a)+F(a)
\end{align}
と書ける．


\newpage

\end{document}
